% Kyle Genova — Curriculum Vitae
% Build: xelatex cv.tex   (output goes to cv.pdf; copy to repo root as CV.pdf)
\documentclass[10.5pt]{article}

\usepackage[letterpaper, margin=0.8in, top=0.6in, bottom=0.65in]{geometry}
\usepackage{fontspec}
\usepackage{xcolor}
\usepackage{enumitem}
\usepackage{titlesec}
\usepackage[hidelinks]{hyperref}
\usepackage{parskip}

\setmainfont{Charter}
\defaultfontfeatures{Ligatures=TeX}

\definecolor{accent}{HTML}{2F7D6E}
\definecolor{body}{HTML}{1E2430}
\definecolor{muted}{HTML}{5C6673}
\color{body}

\hypersetup{colorlinks=true, urlcolor=accent, linkcolor=accent}

\titleformat{\section}{\bfseries\color{body}}{}{0pt}{\large\MakeUppercase}[\vspace{-9pt}\rule{\textwidth}{0.6pt}]
\titlespacing*{\section}{0pt}{9pt}{5pt}

\setlist[itemize]{leftmargin=1.15em, itemsep=1pt, parsep=0pt, topsep=2pt}

\newcommand{\me}{\textbf{Kyle Genova}}
\newcommand{\pub}[3]{\item #1. \textit{#2}. #3}
\newcommand{\entryhead}[3]{\textbf{#1} \hfill {\color{muted}#2}\\ \textit{#3}}

\pagestyle{empty}
\begin{document}

\begin{center}
  {\Huge\bfseries Kyle Genova}\\[5pt]
  {\color{muted} New York, NY \; \textbullet \;
  \href{mailto:kyleagenova@gmail.com}{kyleagenova@gmail.com} \; \textbullet \;
  \href{https://www.kylegenova.com}{kylegenova.com} \; \textbullet \;
  \href{https://scholar.google.com/citations?user=73HIeWcAAAAJ}{Google Scholar} \; \textbullet \;
  \href{https://github.com/kylegenova}{github.com/kylegenova}}
\end{center}

\vspace{-2pt}
\section{Experience}

\entryhead{Google DeepMind — Staff Research Scientist}{Oct.\ 2025 -- Present}{Foundational Research Unit \hfill New York, NY}
\begin{itemize}
  \item Pretraining core contributor to the Gemini, Gemini Omni, and Genie workstreams, focusing on 3D understanding and generation capabilities of frontier models.
  \item Leading the Geo--Gemini Vision project, a 20+ contributor effort that has landed multiple pretraining datasets improving Gemini's spatial understanding.
  \item Contributed multiple petabytes of pretraining data, including the majority of all 3D generation data across Google DeepMind's media-generation models.
\end{itemize}

\entryhead{Google DeepMind — Senior Research Scientist}{May 2024 -- Oct.\ 2025}{Foundational Research Unit \hfill New York, NY}
\begin{itemize}
  \item Contributed pretraining data that made Gemini~3 the first model to surpass human experts at geoguessing.
  \item Led Geo's Foundational Features project, training 3D models to build a petabyte-scale 3D feature database supporting open-vocabulary scene-understanding queries for the Google Maps team.
\end{itemize}

\entryhead{Google Research — Research Scientist}{May 2021 -- May 2024}{Machine Perception --- promoted to Senior Research Scientist, Oct.\ 2023 \hfill Mountain View, CA \& New York, NY}
\begin{itemize}
  \item Published 2D3DNet, the 3D understanding backbone for all of Google Maps since 2021 --- driving over 30 million map updates and supporting routing, HD maps, 3D reconstruction, and indoor Immersive View. Over one quadrillion 3D points served.
\end{itemize}

\entryhead{Google Research — Research Intern (5 internships)}{2016 -- 2021}{New York, Cambridge, and Mountain View}
\begin{itemize}
  \item Research projects in 3D vision, graphics, and machine learning; 9 publications and 4 patents.
\end{itemize}

\section{Education}

\entryhead{Princeton University — Ph.D., Computer Science}{2016 -- 2021}{Advisor: Prof.\ Thomas Funkhouser \hfill Princeton, NJ}
\begin{itemize}
  \item Thesis: \textit{3D Representations for Learning to Reconstruct and Segment Shapes}.
  \item Siebel Scholar (Class of 2021); NSF Graduate Research Fellowship; Gordon Y.S.\ Wu Fellowship in Engineering.
\end{itemize}

\entryhead{Cornell University — B.A., Computer Science}{2012 -- 2016}{GPA 4.17; Phi Beta Kappa (top 3\% of class) \hfill Ithaca, NY}

\section{Selected Publications}

\textbf{24} papers \textbullet\ \textbf{6{,}000+} citations \textbullet\ h-index \textbf{22} \textbullet\ CVPR \textbf{Best Paper Nominee} \textbullet\ {\color{muted} full list on \href{https://scholar.google.com/citations?user=73HIeWcAAAAJ}{Google Scholar}}

\begin{itemize}\small
  \pub{Image Generators are Generalist Vision Learners}{Valentin Gabeur, Shangbang Long, Songyou Peng, Paul Voigtlaender, \ldots, \me, \ldots, Thomas Funkhouser, Jean-Baptiste Alayrac, Radu Soricut (25 authors)}{arXiv, 2026}
  \pub{CityRAG: Stepping Into a City via Spatially-Grounded Video Generation}{Gene Chou, Charles Herrmann, \me, Boyang Deng, Songyou Peng, Bharath Hariharan, Jason Y.\ Zhang, Noah Snavely, Philipp Henzler}{arXiv, 2026}
  \pub{MoMaps: Semantics-Aware Scene Motion Generation with Motion Maps}{Jiahui Lei, \me, George Kopanas, Noah Snavely, Leonidas Guibas}{ICCV 2025}
  \pub{Visual Chronicles: Using Multimodal LLMs to Analyze Massive Collections of Images}{Boyang Deng, Songyou Peng, \me, Gordon Wetzstein, Noah Snavely, Leonidas Guibas, Thomas Funkhouser}{ICCV 2025}
  \pub{SplatTalk: 3D VQA with Gaussian Splatting}{Anh Thai, Songyou Peng, \me, Leonidas Guibas, Thomas Funkhouser}{ICCV 2025}
  \pub{OpenScene: 3D Scene Understanding with Open Vocabularies}{Songyou Peng, \me, Chiyu ``Max'' Jiang, Andrea Tagliasacchi, Marc Pollefeys, Thomas Funkhouser}{CVPR 2023}
  \pub{Panoptic Neural Fields: A Semantic Object-Aware Neural Scene Representation}{Abhijit Kundu, \me, Xiaoqi Yin, Alireza Fathi, Caroline Pantofaru, Leonidas Guibas, Andrea Tagliasacchi, Frank Dellaert, Thomas Funkhouser}{CVPR 2022}
  \pub{Learning 3D Semantic Segmentation with only 2D Image Supervision}{\me, Xiaoqi Yin, Abhijit Kundu, Caroline Pantofaru, Forrester Cole, Avneesh Sud, Brian Brewington, Brian Shucker, Thomas Funkhouser}{3DV 2021 (\textbf{Oral})}
  \pub{IBRNet: Learning Multi-View Image-Based Rendering}{Qianqian Wang, Zhicheng Wang, \me, Pratul P.\ Srinivasan, Howard Zhou, Jonathan T.\ Barron, Ricardo Martin-Brualla, Noah Snavely, Thomas Funkhouser}{CVPR 2021}
  \pub{Local Deep Implicit Functions for 3D Shape}{\me, Forrester Cole, Avneesh Sud, Aaron Sarna, Thomas Funkhouser}{CVPR 2020 (\textbf{Oral})}
  \pub{CvxNet: Learnable Convex Decomposition}{Boyang Deng, \me, Soroosh Yazdani, Sofien Bouaziz, Geoffrey Hinton, Andrea Tagliasacchi}{CVPR 2020 (\textbf{Best Paper Nominee})}
  \pub{Learning Shape Templates with Structured Implicit Functions}{\me, Forrester Cole, Daniel Vlasic, Aaron Sarna, William T.\ Freeman, Thomas Funkhouser}{ICCV 2019}
  \pub{Text-based Editing of Talking-head Video}{Ohad Fried, Ayush Tewari, Michael Zollh\"ofer, Adam Finkelstein, Eli Shechtman, Dan B.\ Goldman, \me, Zeyu Jin, Christian Theobalt, Maneesh Agrawala}{SIGGRAPH 2019}
  \pub{Unsupervised Training for 3D Morphable Model Regression}{\me, Forrester Cole, Aaron Maschinot, Aaron Sarna, Daniel Vlasic, William T.\ Freeman}{CVPR 2018 (\textbf{Spotlight})}
  \pub{An Experimental Evaluation of the Best-of-Many Christofides' Algorithm for the Traveling Salesman Problem}{\me, David P.\ Williamson}{Algorithmica 2017 (\textbf{ESA Selected Publication}); ESA 2015}
\end{itemize}

\section{Patents}

{\color{muted} 6 granted US patents across 4 families}

\begin{itemize}
  \item Neural semantic fields for generalizable semantic segmentation of 3D scenes. \textit{US 12,602,898 (2026)}.
  \item Systems and methods for generating splat-based differentiable two-dimensional renderings. \textit{US 12,288,295 (2025)}.
  \item Convex representation of objects using neural network. \textit{US 11,978,268 (2024); US 11,508,167 (2022)}.
  \item Learning to reconstruct 3D shapes by rendering many 3D views. \textit{US 10,510,180 (2019); US 10,403,031 (2019)}.
\end{itemize}

\section{Academic Service}

\begin{itemize}
  \item \textbf{Area Chair:} CVPR 2026 (\textit{Outstanding Area Chair}); ICCV 2025; SIGGRAPH 2026 (Technical Papers Committee).
  \item \textbf{Committees:} SIGGRAPH Asia 2026 Technical Workshops Committee.
  \item \textbf{Reviewer:} CVPR (2025, 2022, 2021 \textit{Outstanding Reviewer}, 2020, 2019, 2018), ICCV (2023, 2021, 2019), ECCV (2024, 2022, 2020 \textit{Outstanding Reviewer}), NeurIPS (2025 \textit{Top Reviewer}, 2019 \textit{Top Reviewer}), SIGGRAPH (2024, 2023, 2022, 2020), SIGGRAPH Asia (2024, 2023, 2022, 2021, 2019), OpenSUN3D 2024.
\end{itemize}

\section{Honors \& Awards}

\begin{itemize}
  \item Siebel Scholar, Class of 2021 --- The Thomas and Stacey Siebel Foundation, 2020.
  \item NSF Graduate Research Fellowship --- National Science Foundation, 2018--2022.
  \item Graduate Student Teaching Award --- Princeton University, 2018.
  \item Gordon Y.S.\ Wu Fellowship in Engineering --- Princeton University, 2016--2021.
  \item Phi Beta Kappa (top 3\% of class); CRA Outstanding Undergraduate Researcher Nomination --- Cornell University, 2015.
\end{itemize}

\section{Teaching}

\begin{itemize}
  \item \textbf{Princeton University}, Assistant in Instruction: Computer Graphics (COS 426, Spring 2018); Computer Vision (COS 429, Fall 2017). \textit{Graduate Student Teaching Award.}
  \item \textbf{Cornell University}, Teaching Assistant: Analysis of Algorithms (CS 4820); Computer Graphics (CS 4620); Data Structures \& Functional Programming (CS 3110).
\end{itemize}

\section{Skills}

Python, C++, CUDA, JAX, TensorFlow, PyTorch, OpenGL, GLSL.

\end{document}
